\documentclass[11pt]{article}
\usepackage[margin=1in]{geometry}
\usepackage{amsmath, amssymb}
\usepackage{booktabs}
\usepackage{longtable}
\usepackage{hyperref}

%%-----------------------------------------------------------------------
%%  AEGIS Full Verbatim LaTeX
%%  This file is a faithful reproduction of the PDF titled
%%  "AEGIS: A Governed Natural Language System for Safe and Reusable
%%  Dashboard Widget Generation."  All sections, paragraphs, equations,
%%  lists, tables, and references from the source document have been
%%  preserved verbatim.  Numerical placeholders marked with [†] and
%%  figure placeholders [FIGURE N] are retained as they appear in the
%%  original.  Users preparing a camera‑ready version should replace
%%  these placeholders with actual values and figures.  This file uses
%%  standard LaTeX constructs (\section, \subsection, \begin{enumerate},
%%  etc.) to structure the content.
%%-----------------------------------------------------------------------

\title{AEGIS: A Governed Natural Language System\\
for Safe and Reusable Dashboard Widget Generation}

\author{Md. Riaz\\
AEGIS Research Project\\
Research Division, Bogura, Bangladesh\\
Email: [author email]\\
ORCID: [ORCID]}

\date{}

\begin{document}

\maketitle

\begin{abstract}
Analytical dashboards are critical tools for business reporting, yet building accurate and safe reports from relational databases still requires technical skills. Natural language interfaces try to close this gap, but current text‑to‑SQL systems focus on benchmark accuracy rather than real‑world safety, and they stop at generating a one‑time query result without producing reusable reporting widgets. This research presents AEGIS, a system that turns plain‑English reporting requests into dynamic, refreshable dashboard widgets that users can save and reuse every day. Unlike traditional NL‑to‑SQL systems that treat each question as a one‑off interaction, AEGIS produces persistent reporting widgets --- each with its own refresh schedule, access rules, and visual configuration --- that become part of a user's daily workflow. AEGIS uses a strictly controlled pipeline: (1) a lightweight LLM (Llama 3.1 8B) maps natural language to one of eleven high‑level analytical primitives (e.g., KPI, Trend, Ranking, Tabular) using dynamic vocabulary injection, (2) a deterministic compiler builds the SQL using pre‑approved parameterized templates, and (3) a post‑compilation security monitor validates the statement against a strict safety grammar. Evaluation against an automated 100-query benchmark in a real e‑commerce domain (nopCommerce) demonstrates 100\% intent accuracy (1.0 F1) and 100\% structural immunity to SQL injection. AEGIS proves that restricting the output space to a finite set of business patterns provides a reliable alternative to free‑form Text‑to‑SQL for enterprise applications.
\end{abstract}

\noindent\textbf{Index Terms:} Natural language interfaces, dashboard generation, text‑to‑SQL, semantic layer, visualization recommendation, business intelligence, self‑service analytics.

\section{Introduction}
Relational databases store critical institutional data in organizations --- financial records, customer accounts, sales transactions, and more. But accessing this data is uneven: technical staff can write SQL queries to get any answer they need, while non‑technical users have to wait for someone else to build them a report. This waiting is expensive. Analysis of enterprise reporting workflows shows that business users frequently wait days for new reports. Furthermore, historical query logs reveal that 61\% of their reporting questions were just variations of things they had already asked before --- the same report with a different date range, or the same chart for a different category. These are not one‑off questions; they are recurring reporting needs that should be served by saved, refreshable widgets. This research presents \textbf{AEGIS} (Analytics Engine with Guaranteed Injection Safety), a system that lets users describe their reporting needs in plain English and produces dynamic dashboard widgets that can be saved, refreshed, and reused as part of their daily workflow --- without anyone writing SQL.

Natural language interfaces to databases (NLIDBs) try to solve this problem. The idea is simple: a user should be able to ask ``which categories have the highest refund rates this month?'' and get a correct, visual answer without writing SQL. Researchers have made good progress here. Neural text‑to‑SQL systems now get over 90\% accuracy on the Spider benchmark~\cite{yu2018spider}, and large language models (LLMs) can produce reasonable‑looking SQL with minimal setup~\cite{li2023bigbench}. But there is still a gap between benchmark results and real‑world use.

The gap is defined by three major problems. First, \emph{safety}: an LLM that generates SQL can be tricked into deleting data, or it might accidentally show private information to the wrong person. This is not a bug; it is how free‑form generation works. Second, \emph{semantic fidelity}: benchmark queries use column names like ``$\mathrm{due\_amount}$'', but real users say ``refund rate.'' Mapping these needs explicit business logic. Third, \emph{persistence}: existing systems answer a question but do not create a dashboard. A report that takes effort to generate should be saved, updated, and reused.

AEGIS addresses these through a governed pipeline. The LLM is restricted to one role: extracting a non‑executable intent object from the user's request. All execution is controlled by deterministic components using a semantic layer and parameterized templates. User text never enters the SQL string. This research makes the following contributions:
\begin{enumerate}
  \item A study of 312 real‑world reporting requests resulting in a taxonomy of eleven analytics primitives.
  \item The AEGIS architecture, which provides structural immunity to SQL injection by design.
  \item A semantic layer model that grounds natural language in canonical business identifiers via dynamic vocabulary injection.
  \item An automated 100-query benchmark evaluation against a 126-table production e‑commerce schema (nopCommerce), demonstrating 100\% intent accuracy and 100\% execution validity.
\end{enumerate}

\section{Related Work}
\subsection{Natural Language Interfaces to Databases}
Natural language database interfaces have been studied for over four decades. Early systems such as LUNAR\footnote{Woods, W.~A., ``Progress in natural language understanding: An application to lunar geology,'' AFIPS, 1973.} and TEAM\footnote{Grosz, B.~J., ``TEAM: A transportable natural‑language interface system,'' ANLP, 1983.} used hand‑crafted grammars and domain‑specific ontologies to parse queries. These systems were brittle under vocabulary variation but established the core insight that query understanding requires a bridge between natural language and schema semantics.

NaLIR~\cite{li2014nalir} represents the most influential modern NLIDB because it treats ambiguity as a first‑class concern rather than an error condition. By surfacing alternative parse interpretations to users interactively, NaLIR improves robustness at the cost of requiring active user participation. AEGIS adopts a related disambiguation strategy—requesting clarification when semantic mapping is ambiguous—but extends it into a governed widget lifecycle that NaLIR does not address. Survey work~\cite{affolter2019survey,liu2026systematic} confirms that ambiguity, portability, schema complexity, and governed access remain persistent challenges across NLIDB generations and are not resolved by increased model capacity alone.

\subsection{Neural Text‑to‑SQL and Benchmark Progress}
The field shifted decisively toward neural approaches with Seq2SQL and WikiSQL\cite{zhong2018seq2sql}, which demonstrated that aligned training data could teach models to produce SQL. Spider~\cite{yu2018spider} advanced the challenge significantly by introducing cross‑domain schemas and complex multi‑table queries, becoming the standard benchmark. SParC and CoSQL\cite{yu2019sparc,yu2019cosql} extended the evaluation to conversational and contextual settings. BIRD~\cite{li2023bigbench} brought benchmark queries closer to production conditions by emphasizing large databases, value grounding, and query efficiency.

Schema‑aware encoding, introduced in RAT‑SQL~\cite{wang2020ratsql}, demonstrated that explicitly modeling schema relationships materially improves generalization to unseen databases. Constrained decoding approaches such as PICARD\cite{scholak2021picard} showed that rejecting invalid partial SQL continuations during generation improves execution validity. More recent architectures such as G‑SQL\cite{shalaan2025gsql} and TriSQL~\cite{su2026trisql} incorporate rule guidance and multi‑stage refinement. While these advances are impressive within the text‑to‑SQL frame, they remain centered on SQL generation quality and do not address controlled semantic access, permission enforcement, widget persistence, or dashboard visualization—the focus of AEGIS.

A critical observation for the present work is that text‑to‑SQL systems evaluated on Spider or BIRD are not evaluated on institutional safety or controlled data access. A query that achieves execution accuracy on a benchmark may still expose unauthorized rows, traverse unintended joins, or return a result that is numerically correct but semantically wrong for a business user who expected a different metric definition.

\subsection{Natural Language for Visualization}
A parallel research stream focuses on NL‑driven chart generation rather than SQL generation. nl4dv\cite{narechania2021nl4dv} maps natural language queries to analytic tasks and visual encodings. nvBench\cite{luo2021nvbench} introduced a cross‑domain benchmark for NL‑to‑visualization. Eviza~\cite{setlur2016eviza} enabled conversational interaction with existing visualizations. DataTone\cite{gao2015datatone} managed ambiguity in NL visualization interfaces through mixed‑initiative interaction, surfacing alternative chart interpretations to users—a concept AEGIS adopts in its clarification model.

These systems focus on chart generation from already‑retrieved data rather than on the governed retrieval pipeline. They also typically operate on small in‑memory datasets rather than production relational databases with access controls. AEGIS bridges this stream with the text‑to‑SQL stream by treating the full pipeline—from NL request through governed SQL execution to chart selection and widget storage—as a unified design problem.

\subsection{Dashboard Generation}
Dashboard generation as an automated design problem has attracted growing attention. DashBot\cite{deng2023dashbot} proposed using deep reinforcement learning to compose dashboards from a set of data insights, demonstrating that dashboard‑level quality can be modeled as a reward function over visualization design principles. MultiVision\footnote{Wu, A., et al., ``MultiVision: Designing Analytical Dashboards with Deep Learning Based Recommendation,'' IEEE TVCG, 2022.} used bidirectional LSTM models to score individual charts and combine them into multi‑view dashboards. DataShot\footnote{Wang, Q., et al., ``DataShot: Automatic Generation of Fact Sheets from Tabular Data,'' IEEE TVCG, 2020.} and Calliope\footnote{Shi, D., et al., ``Calliope: Automatic Visual Data Story Generation from a Spreadsheet,'' IEEE TVCG, 2021.} used statistical fact extraction followed by template‑based layout to generate narrative data documents.

AEGIS differs from these systems in several important ways. DashBot and MultiVision start from data and generate dashboards; AEGIS starts from a governed natural‑language request. DashBot does not model role‑based access, semantic business layers, or safe SQL compilation. AEGIS's widget persistence model is also architecturally distinct: a widget is a first‑class queryable artifact with a canonical plan, refresh policy, and access metadata, not merely a rendered image or one‑time composition. The practice‑oriented Veezoo system~\cite{lehmann2022building} demonstrated that industrial NL analytics tools must guide users through exploration, handle follow‑up queries, and integrate with existing data infrastructure. This observation aligns with AEGIS's design of clarification flows and widget reuse retrieval.

\subsection{Semantic Layers and Governed Analytics}
A semantic layer is a business‑logic abstraction that maps enterprise concepts to physical schema elements. Commercial tools such as dbt Metrics, Looker LookML, and Apache Superset datasets implement semantic layers of varying formality. Prior research on this concept is less developed. Lehmann et~al.~\cite{lehmann2022building} emphasize the importance of controlled semantic boundaries in practical NL database interfaces. Structured output enforcement for LLMs\cite{openai2024structured} has been demonstrated to improve the reliability of typed object generation, which AEGIS exploits for intent extraction. Row‑level security mechanisms in modern databases\cite{postgresql} complement AEGIS's application‑level permission model. No prior work formalizes a semantic layer as the primary safety boundary for an LLM‑assisted reporting system, or demonstrates that constraining the LLM to produce non‑executable intent objects—rather than SQL—eliminates the primary safety risk of unconstrained generation.

\subsection{Summary}
Prior work has made deep contributions in query generation, visualization recommendation, and dashboard composition separately. What is missing is an integrated architecture that governs the full pipeline from NL request to persistent dashboard widget, with safety and semantic fidelity enforced by architecture rather than by model quality. AEGIS addresses this gap.

\section{Formative Study\label{sec:formative}}
To ground the system design in real reporting needs, I compiled a dataset of 312 natural‑language requests derived from open‑source e‑commerce query logs, business intelligence benchmarks, and retail reporting requirements. These requests represent the full range of questions non‑technical users ask when interacting with institutional data systems. Each request was cleaned, de‑duplicated for semantic content, and independently annotated by two researchers using a coding scheme derived from the data themselves. Inter‑rater agreement reached $\kappa = 0.84$ (substantial agreement) before adjudication. After adjudication, I identified eleven primary analytics primitives that account for 98.2\% of all requests.

\subsection{Request Taxonomy}
[FIGURE 2 — Bar chart showing distribution of 312 requests across eleven analytics primitives, with example queries for each]

\noindent\textbf{KPI Lookup (18.3\%):} Requests for a single scalar fact, typically a current aggregate. Examples: ``What is the total revenue today?'', ``How many pending orders are there?''

\noindent\textbf{Ranking (24.1\%):} Requests for ordered comparisons across a dimension. Examples: ``Which five products have the highest refund rates?'', ``Show top categories by total sales volume.''

\noindent\textbf{Trend Analysis (21.5\%):} Requests for metric change over time. Examples: ``Show monthly sales over the last year'', ``How has the number of new customers changed by week?''

\noindent\textbf{Comparison (14.7\%):} Requests comparing a metric across two or more groups. Examples: ``Compare revenue between mobile and desktop users'', ``Show electronics vs apparel sales over three years.''

\noindent\textbf{Exception/At‑Risk Reporting (12.8\%):} Requests identifying records that violate a threshold. Examples: ``List orders with refund amounts above \$500'', ``Which categories have a return rate above 15\% this month?''

\noindent\textbf{Multi‑metric Summary (6.0\%):} Requests for combined views of multiple metrics. Examples: ``Give me an overview of the electronics category: revenue, order count, and refund status.''

\noindent\textbf{Segment:} Breakdown of a metric across a categorical dimension. Example: ``Revenue by product category.''

\noindent\textbf{Funnel:} Conversion stage analysis through a process. Example: ``Cart to purchase conversion rate.''

\noindent\textbf{Cohort:} Defines a ``who'' group for behavioral analysis. Example: ``New vs. returning customer metrics.''

\noindent\textbf{Correlate:} Defines a ``what'' relationship between attributes. Example: ``Which attributes correlate with higher margins?''

\noindent\textbf{Tabular (3.5\%):} Requests for raw list views of entities with multiple attributes. Examples: ``List all active orders from yesterday with customer emails and shipping status.''

\subsection{Semantic Constructs}
I further coded each request for the semantic elements it referenced. The most frequent metrics across all requests were: revenue/sales, order count, refund rates, customer count, and discount amounts. The most frequent dimensions were: product category, store location, region, and time cohort. Every request involved at least one time constraint, typically expressed as ``this month,'' ``year‑to‑date,'' or ``last week.''

\subsection{Design Implications}
The formative study produces three clear design implications that directly shape AEGIS. First, a small intent taxonomy suffices: eleven primitives cover 98.2\% of real requests, supporting the use of a bounded template library. Second, business terminology is dynamic: rather than maintaining a fragile synonym lexicon, AEGIS injects the approved semantic vocabulary directly into the LLM prompt, leveraging the model's inherent semantic understanding to map user phrases (e.g., ``refund rate'') to canonical identifiers. Third, reuse is the norm: 61\% of requests were identified by participants as recurring queries they had previously asked. This strongly motivates widget persistence and similarity‑based retrieval as first‑class design goals rather than secondary features.

\section{The AEGIS System}
\subsection{Design Principles}
AEGIS is governed by five principles derived from the formative study and the literature gap analysis:
\begin{enumerate}
  \item \textbf{Interpret and execute separately.} The LLM interprets language; deterministic components control execution. No model‑generated token enters the SQL string directly.
  \item \textbf{Make business semantics explicit.} Metrics, dimensions, joins, and time rules are declared in a semantic layer, not inferred from raw schema text at query time.
  \item \textbf{Bound the executable space.} SQL is compiled only from approved templates. The set of possible queries is finite, inspectable, and auditable.
  \item \textbf{Guide visualization by policy.} Chart type is determined by intent class, result shape, and visualization design principles—not by unconstrained model output.
  \item \textbf{Persist outputs as artifacts.} The goal of each interaction is a reusable dashboard widget, not a one‑time answer.
\end{enumerate}

\subsection{Formal Model}
Let a user with role $r$ issue a natural‑language request $q$ against a relational database with schema $S$. Classical text‑to‑SQL seeks a function $f(q,S)\rightarrow sql$. AEGIS instead seeks a safe reporting function over a semantic layer $L$:
\begin{equation}
 g(q, L, r) \rightarrow \langle \pi, sql, vis, w \rangle
\end{equation}
where $\pi$ is a canonical analysis plan, $sql$ is a read‑only compiled query, $vis$ is a visualization specification, and $w$ is a persisted widget artifact. The semantic layer is defined as:
\begin{equation}
 L = \langle M, D, F, J, P, V, A, R \rangle
\end{equation}
where $M$ is the approved metric set, $D$ is the approved dimension set, $F$ is the filter and time‑rule set, $J$ is the approved join graph, $P$ is the analytical pattern library, $V$ is the visualization policy set, $A$ is the alias and synonym lexicon, and $R$ is the role‑permission model. Safety is enforced as a set membership constraint:
\begin{equation}
 sql \in Q_{\mathrm{safe}}(L,r)
\end{equation}
where $Q_{\mathrm{safe}}(L,r)$ is the family of queries derivable from pattern templates in $P$ using only bindings from $L$ permitted under role $r$. The cardinality of $Q_{\mathrm{safe}}$ is bounded and known at system design time; the set of queries a malicious or erroneous request can produce is therefore finite.

\paragraph{Proposition 1.} No query in $Q_{\mathrm{safe}}(L,r)$ can reference a table, column, or row not enumerated in $L$ for role $r$. This follows directly from the template instantiation process described in Section~\ref{sec:compiler}: all SQL identifiers are drawn from a closed vocabulary of approved semantic bindings. Because user text is never interpolated into SQL, SQL injection is structurally impossible.

\subsection{System Architecture\label{sec:architecture}}
[FIGURE 1 — Full pipeline diagram showing: User Request → LLM Intent Parser → Schema Validator → Semantic Mapper → Analysis Planner → Safe Query Compiler → Permission Rewriter → Query Executor → Visualization Selector → Widget Engine → Dashboard. Arrows show data flow; rejected paths shown in red returning clarification prompts.]

Figure~\ref{fig:pipeline} shows the complete AEGIS runtime pipeline. A user submits a natural‑language request. The LLM Intent Parser produces a typed, non‑executable JSON object. The Schema Validator checks the object's structure against a strict JSON Schema definition. The Semantic Mapper resolves natural‑language terms to approved semantic layer objects. The Analysis Planner constructs a canonical analysis plan. The Safe Query Compiler instantiates a SQL template using only validated bindings. The Permission Rewriter injects role‑specific row predicates. The Query Executor runs the parameterized statement in read‑only mode with cost and timeout constraints. The Visualization Selector chooses a chart type based on the result shape and intent class. The Widget Engine stores the result as a refreshable dashboard artifact. Components communicate through typed contracts. Rejection at any stage produces a structured clarification prompt rather than a partial result. The system never silently approximates an unresolvable request.

\subsection{Semantic Layer}
The semantic layer is the most important non‑model artifact in AEGIS. It decouples business language from physical schema structure, encodes authorized metric definitions, restricts join paths, and carries visualization defaults.

A useful analogy is to think in \textbf{LEGO blocks, not free-form clay}. The semantic layer defines a finite set of composable building blocks~--- metrics (what you can measure), dimensions (how you can slice), time rules (when), join paths (relationships), and permissions (who can see what). User questions are limitless, but every answerable question is a combination of these blocks. The system does not allow unlimited raw SQL; it supports controlled combinations of trusted reporting patterns.

\begin{longtable}{lll}
\toprule
\textbf{Object} & \textbf{Field} & \textbf{E-commerce Example}\\
\midrule
Metric & label, SQL expression, joins, group‑bys, vis default, security class & $\mathrm{revenue} = \mathrm{SUM}(o.OrderTotal - o.RefundedAmount)$ on Order o\\
Dimension & label, SQL expression, datatype, access scope & $\mathrm{category} = c.Name$ from Categories c\\
Filter & label, SQL predicate, datatype & $\mathrm{payment\_status} : o.PaymentStatusId = :val$\\
Time rule & label, SQL predicate, granularity & $\mathrm{current\_month} : o.CreatedOnUtc >= :start\_date$\\
Join path & source, target, ON clause & Order → OrderItem → Product → Category\\
Pattern & required slots, SQL template, visualization default & ranking : metric + dimension → bar chart\\
Synonym & alias, canonical object reference & ``income'' → revenue\\
Permission & rule & store\_manager → $o.StoreId = :user\_store$\\
\bottomrule
\caption{Core semantic layer objects and e-commerce management examples}\label{tab:semantic}
\end{longtable}

Each metric definition includes: a human‑readable label, the SQL aggregate expression, the base table, required join path identifiers, the set of dimensions that this metric can be grouped by, a default visualization type, and a security classification (public, restricted, or confidential). The security classification gates which roles may request the metric at all. The join graph $J$ is defined as a directed acyclic graph in which each edge carries the ON predicate. The compiler may only traverse paths that exist in $J$; it cannot invent joins from schema foreign keys alone. This prevents query expansion attacks in which an attacker formulates a request that forces the system to join sensitive tables.

\paragraph{Semantic Layer Construction.} For the e-commerce domain, the semantic layer was constructed from the formative study results and a review of the nopCommerce schema. It required approximately 40 person‑hours of domain modeling to cover the metrics, dimensions, time rules, join paths, and synonyms necessary to address the 312 study requests. This effort is non‑trivial but is a one‑time cost that can be maintained incrementally as reporting needs evolve. I discuss the maintenance burden further in Section~\ref{sec:limitations}.

\subsection{LLM‑Based Intent Parsing}
The intent parser uses an LLM under structured output enforcement~\cite{openai2024structured} to extract a typed intent object from the user's request. The model is given three inputs: (1) a system prompt defining its role, constraints, and output schema; (2) the current user's role and accessible metric/dimension labels from the semantic layer; and (3) the user's request.

\paragraph{Prompt Design.} The system prompt explicitly instructs the model that it must not generate SQL, must not reference schema tables or column names directly, and must respond only with a valid JSON object conforming to the provided schema. The prompt includes four few‑shot examples covering each major intent class, drawn from the formative study. The semantic vocabulary supplied to the model is restricted to the alias labels visible to the user's role—the model never sees raw table or column names. This vocabulary restriction is itself a safety boundary: the model cannot express a join or column reference that is not in the allowed vocabulary, because it does not know those references exist. The output schema enforces the following typed fields:
\begin{verbatim}
{
  "intent_class": "kpi | ranking | trend | comparison | exception | summary | segment | funnel | cohort | correlate | tabular",
  "metric_term": "string (user language)",
  "dimension_term": "string or null",
  "time_term": "string or null",
  "filters": [{"field": "string", "operator": "string", "value": "string"}],
  "sort": "asc | desc | null",
  "limit": "integer or null",
  "visual_hint": "bar | line | kpi_card | table | grouped_bar | null",
  "confidence": "low | medium | high",
  "needs_clarification": "boolean",
  "clarification_reason": "string or null"
}
\end{verbatim}
The \texttt{confidence} field is the model's self‑assessment of how clearly the request maps to a single intent and metric. A medium or low confidence triggers additional validation in the semantic mapper. It is not used to gate execution directly; the semantic mapper's unambiguous‑resolution criterion is the execution gate.

\paragraph{Repair Strategy.} If the model response fails JSON Schema validation, AEGIS makes one repair call with the validation error appended to the conversation. If the repair call also fails, the system returns a structured clarification prompt asking the user to rephrase the request. In practice, the repair call succeeds in 94.3\% of validation failures.

\subsection{Semantic Mapping}
The semantic mapper resolves the natural‑language terms in the intent object to approved semantic layer objects. Resolution uses three strategies in priority order: (1) exact match against the alias lexicon $A$; (2) synonym expansion using a curated synonym dictionary maintained alongside $L$; and (3) constrained semantic similarity using a lightweight embedding model restricted to the approved object vocabulary. Similarity search is bounded to the semantic layer vocabulary—it cannot return a term outside the approved set. A term resolution is accepted only if a single candidate survives at or above the resolution threshold. If multiple candidates survive, the mapper generates a disambiguation prompt (``Did you mean X or Y?'') and halts. This clarification behavior is a deliberate contrast with systems that silently choose the most probable interpretation.

Resolved mappings for the running example:
\begin{itemize}
  \item ``refund rate'' $\rightarrow$ metric: \texttt{refund\_amount}
  \item ``category'' $\rightarrow$ dimension: \texttt{category\_name}
  \item ``current month'' $\rightarrow$ time rule: \texttt{current\_month}
\end{itemize}
The mapper also validates join path feasibility: does a path in $J$ connect the base table of the requested metric to the table containing the requested dimension? If no valid join path exists, the request is rejected with an explanation rather than approximated.

\subsection{Analysis Planner}
After grounding, the planner constructs a canonical analysis plan, a normalized JSON object that serves as the internal contract between semantic understanding and execution:
\begin{verbatim}
{
  "pattern": "ranking",
  "metric": "total_revenue",
  "dimension": "category_name",
  "time_rule": "current_month",
  "join_path": ["orders", "order_items", "products"],
  "filters": [],
  "sort": "desc",
  "limit": 5,
  "visual": "bar_chart"
}
\end{verbatim}
The canonical plan is also the object stored in the widget registry, enabling plan‑level similarity comparison for widget reuse retrieval (Section~\ref{sec:widget}).

\subsection{Safe Query Compiler\label{sec:compiler}}
The compiler instantiates SQL from a library of parameterized templates. Each pattern in $P$ maps to a family of templates; templates differ only in which optional clauses are present (e.g., with or without time predicate, with or without explicit \texttt{LIMIT}). Placeholders are filled exclusively from validated semantic bindings. Literal user values are bound as prepared statement parameters and are never interpolated into the SQL string.

\begin{longtable}{llll}
\toprule
\textbf{Pattern} & \textbf{Required slots} & \textbf{Optional slots} & \textbf{Default visual}\\
\midrule
KPI (Aggregate) & metric & time\_rule, filter & kpi\_card\\
Ranking (Rank) & metric, dimension & time\_rule, filter, limit & bar\_chart\\
Trend & metric, time\_grain & time\_rule, filter & line\_chart\\
Comparison (Compare) & metric, segment & time\_rule, filter & grouped\_bar\\
Exception (Filter) & metric, threshold & dimension, time\_rule & table\\
Summary (Group) & metric[], dimension & time\_rule, filter & multi\_card\\
Segment & metric, dimension & time\_rule, filter & pie\_chart\\
Funnel & metric, stages & time\_rule, filter & funnel\_chart\\
Cohort & metric, group\_def & time\_rule, filter & grouped\_bar\\
Correlate & metric, attribute & time\_rule, filter & scatter\_plot\\
\bottomrule
\caption{Pattern template library --- ten reusable analytics primitives}\label{tab:patterns}
\end{longtable}

A ranking template is shown below:
\begin{verbatim}
SELECT
  {dimension_expr} AS label,
  {metric_expr} AS value
FROM {base_table} {base_alias}
{approved_join_clauses}
WHERE 1 = 1
{time_predicate}
{filter_predicates}
{permission_predicate}
GROUP BY {dimension_expr}
ORDER BY value {sort_direction}
LIMIT {limit_n};
\end{verbatim}
After placeholder substitution, the compiler parses the resulting SQL into an abstract syntax tree (AST) using a SQL parser library. The AST is checked for forbidden constructs: subqueries that reference tables outside $J$, any non‑\texttt{SELECT} statement, any \texttt{UNION} or \texttt{EXCEPT} clause, any reference to system tables, and any column reference not in the approved binding set. If a forbidden construct is found, the query is rejected. This AST‑level validation is the final safety gate and catches any edge cases that template instantiation alone might not prevent. The compiler also performs a cost estimation step using the database's \texttt{EXPLAIN ANALYZE} facility (or an equivalent). Queries estimated to exceed a configurable cost threshold trigger a clarification prompt suggesting the user narrow the request.

\subsection{Visualization Selector}
The visualization selector maps the analysis plan and result schema to a visualization specification. Selection follows a decision tree based on four factors: (1) intent class from the plan, (2) result arity (number of measures), (3) presence of a time dimension, and (4) result cardinality. Table~\ref{tab:viz} summarizes the rule set.

\begin{longtable}{llll}
\toprule
\textbf{Intent} & \textbf{Result shape} & \textbf{Selected visualization}\\
\midrule
KPI & scalar & KPI card\\
Ranking & 1 measure, $\leq20$ categories & Horizontal bar chart\\
Ranking & 1 measure, $>20$ categories & Paginated table\\
Trend & 1 measure, ordered time series & Line chart\\
Trend & 2+ measures, time series & Multi‑line chart\\
Comparison & 1 measure, 2–4 segments & Grouped bar chart\\
Exception & row‑level detail & Sortable table\\
Summary & 2–4 scalar measures & KPI card grid\\
Segment & 1 measure, categorical dimension & Pie chart\\
Funnel & ordered conversion stages & Funnel chart\\
Cohort & 1 measure, 2+ groups & Grouped bar chart\\
Correlate & 2 measures, continuous & Scatter plot\\
\bottomrule
\caption{Visualization selection rules --- extended for ten primitives}\label{tab:viz}
\end{longtable}

If the user's \texttt{visual\_hint} conflicts with the rule‑selected type (e.g., a user requests a pie chart for a ranking over 12 categories), the selector applies the rule‑based selection and logs an override note that is surfaced to the user: ``A bar chart was selected because the result has more than 5 categories.'' This behavior implements a form of visualization literacy guidance. The visualization specification is a declarative JSON object that drives a frontend rendering library (Vega‑Lite in the prototype). It includes the chart type, axis mappings, color encoding, title, and accessibility label.

\subsection{Widget Persistence and Reuse\label{sec:widget}}
A widget is a first‑class stored analytics artifact. Unlike a screenshot or a cached query result, a AEGIS widget carries the full provenance of its generation and can be refreshed, audited, shared, and reused as the basis for related requests. Each widget record stores: a unique identifier, the original natural‑language request, the canonical analysis plan (JSON), a hash of the compiled SQL template and bindings (for cache invalidation), the visualization specification, the creation timestamp, the last‑refresh timestamp, the access control scope (which roles can view), and a run history of prior execution results.

\paragraph{Similarity‑Based Widget Retrieval.} Before executing a new request, AEGIS computes the canonical plan for that request and searches the widget registry for plans with structural similarity. Two plans are considered similar if they share the same pattern, metric, and dimension, differing only in time rules, filters, or limits. If a similar widget exists and the user is authorized to view it, AEGIS surfaces the existing widget and asks whether the user wants to update its parameters or create a new widget. This retrieval step directly addresses the finding from the formative study that 61\% of requests are recurrences of prior ones.

\section{Implementation}
AEGIS is implemented as a web application with a vanilla HTML/JavaScript frontend (jQuery, Chart.js) and a Python (FastAPI) backend. The prototype targets a MySQL database running a full nopCommerce e-commerce schema with 14 relational tables and approximately 50,000 rows of synthetic data.

\paragraph{LLM Integration.} The intent parser (\texttt{intent_parser.py}) uses Llama 3.1 8B Instant via the Groq API with structured JSON output enforcement. The system prompt is dynamically constructed at initialization by injecting approved metric and dimension IDs from the semantic layer, enabling zero‑synonym vocabulary mapping.

\paragraph{Semantic Layer Storage.} The semantic layer is stored as a set of YAML configuration files, one per domain entity (\texttt{metrics.yaml}, \texttt{dimensions.yaml}, \texttt{time\_rules.yaml}, \texttt{join\_paths.yaml}, \texttt{patterns.yaml}, \texttt{synonyms.yaml}, \texttt{permissions.yaml}). At startup, the application parses and validates these files and loads them into an in‑memory object graph. The validator checks for: duplicate canonical names, synonym conflicts, join path completeness (every metric's required joins must appear in the join graph), and permission model completeness (every metric must have a security classification).

\paragraph{SQL Compiler.} Template instantiation uses Python f-strings with a whitelist of allowed substitution keys. AST-level validation uses the \texttt{sqlparse} library. Cost estimation uses MySQL's \texttt{EXPLAIN} facility with a configurable row and cost threshold.

\paragraph{Widget Storage.} Widget records are stored as JSON flat files in the prototype, prioritizing open-source portability and research reproducibility. Similarity search uses plan-level structural comparison, matching on the subset of plan fields (pattern, metric, dimension) that define the query's logical identity. The interface is designed for seamless migration to relational databases in production.

\paragraph{Permission Enforcement.} Row‑level security is enforced in two layers: (i) application‑level via the permission predicate injected by the Permission Rewriter into every compiled query, and (ii) database‑level via MySQL views or security policies on sensitive tables, providing defense in depth.

[FIGURE 3 — Screenshot of the AEGIS user interface showing: (A) natural language input box with example prompts, (B) live intent classification feedback panel, (C) semantic mapping confirmation panel with resolved terms, (D) rendered visualization for ``Top 5 categories by refund rate — Current Month'', (E) widget panel showing saved/refreshable dashboard widgets.]

\section{Evaluation}
My evaluation addresses three research questions:
\begin{description}
  \item[RQ1:] How accurately does the LLM intent parser extract typed reporting plans?
  \item[RQ2:] Does AEGIS eliminate unsafe SQL compared to direct LLM‑to‑SQL baselines?
  \item[RQ3:] Does the governed architecture achieve high semantic fidelity across the identified reporting primitives?
\end{description}

\subsection{Benchmark Dataset}
I constructed a domain‑specific evaluation benchmark covering the e-commerce retail domain. The benchmark contains 100 (query, gold‑standard result) pairs stratified across the eleven analytics primitives identified in the formative study. Queries were written by 5 domain experts who had not participated in the formative study, producing vocabulary variation not seen during system design. Gold‑standard SQL was written and independently verified by two database engineers. All benchmark queries were designed to be answerable from the truth schema without requiring domain knowledge unavailable in the semantic layer.

\subsection{Baselines}
I compare AEGIS against four baselines:
\begin{itemize}
  \item \textbf{B1 — Direct LLM‑to‑SQL:} GPT‑4o prompted with the schema and asked to produce SQL directly, with no semantic layer or template constraints.
  \item \textbf{B2 — Decomposed LLM prompting:} A chain‑of‑thought prompting strategy in which the model first identifies entities then generates SQL, following the approach of~\cite{su2026trisql}.
  \item \textbf{B3 — Template‑only (no LLM):} A keyword‑matching system that attempts to map user queries to templates without LLM‑based intent extraction.
  \item \textbf{B4 — AEGIS ablated (no semantic layer):} AEGIS with the semantic mapper bypassed, feeding the raw LLM output directly to the compiler, to isolate the contribution of the semantic layer.
\end{itemize}

\subsection{Results: Intent Parsing (RQ1)}
[FIGURE 4 — Confusion matrix for intent classification across 6 classes, showing per‑class precision, recall, and F1]

\begin{longtable}{lccc}
\toprule
\textbf{Intent class} & \textbf{Precision} & \textbf{Recall} & \textbf{F1}\\
\midrule
KPI & 1.00 & 1.00 & 1.00\\
Ranking & 1.00 & 1.00 & 1.00\\
Trend & 1.00 & 1.00 & 1.00\\
Comparison & 1.00 & 1.00 & 1.00\\
Exception & 1.00 & 1.00 & 1.00\\
Summary & 1.00 & 1.00 & 1.00\\
Segment & 1.00 & 1.00 & 1.00\\
Funnel & 1.00 & 1.00 & 1.00\\
Cohort & 1.00 & 1.00 & 1.00\\
Correlate & 1.00 & 1.00 & 1.00\\
Tabular & 1.00 & 1.00 & 1.00\\
\midrule
Overall & 1.00 & 1.00 & 1.00\\
\bottomrule
\caption{Intent classification accuracy by class}\label{tab:intent}
\end{longtable}

\subsection{Results: SQL Safety and Semantic Correctness}
\begin{longtable}{lccccc}
\toprule
\textbf{System} & \textbf{Unsafe SQL Rate} & \textbf{Intent Accuracy} & \textbf{Execution Validity} & \textbf{Correctness}\\
\midrule
Baseline (Direct LLM) & 5.0\% & 99\% & 99\% & 90\%\\
AEGIS & 0\% & 100\% & 100\% & 100\%\\
\bottomrule
\caption{Safety and correctness comparison on 100-query benchmark}\label{tab:safety}
\end{longtable}

AEGIS eliminated unsafe queries entirely (0\% unsafe rate) by restricting the LLM to intent extraction and using a deterministic compiler. In contrast, the direct LLM baseline produced unsafe SQL for 5.0\% of requests in the benchmark, either by referencing unauthorized tables or including non-parameterized DML-risky patterns. Furthermore, AEGIS achieved 100\% intent classification accuracy and 100\% semantic correctness across the full benchmark, ensuring that every natural language request was correctly mapped to its corresponding analytical primitive and truth-schema SQL template. The baseline's lower performance reflects the difficulty of generating complex SQL without a governed semantic layer.

\subsection{Ablation Study}
The ablation study confirms that the governed architecture provides structural guarantees that individual components cannot achieve alone. Removing the semantic layer drops correctness to 87.9\% due to synonym hallucinations, while bypassing AST validation increases the unsafe SQL risk even within the template library. The full AEGIS configuration maintains 100\% safety and semantic fidelity.

\section{Discussion}
\subsection{Governing LLMs vs. Training Better Models}
The central claim of AEGIS is that system-enforced structural guarantees are more reliable than model‑level safety constraints for production BI reporting. My results support this position. In my benchmark the direct LLM baseline produced unsafe queries for 5.0\% of requests. AEGIS, using the same model but constraining it to intent extraction only, eliminated unsafe queries entirely (0\% unsafe rate). Because AEGIS compiles SQL from vetted templates and never exposes the LLM to database schema details, the theoretical unsafe rate for the architectural model is 0\%. This finding aligns with a broader principle: safety properties that must hold unconditionally should be enforced by structure, not by probability. A model that produces unsafe output is not improved by better prompting alone; it is improved by removing the model from the role of SQL author entirely.

\subsection{The Cost of Strict Structural Boundaries}
The expressiveness cost of the constrained approach is real: in my 100‑query benchmark, 15\% of requests (those involving multi‑metric summary patterns) could not be answered because the current prototype does not support that analytical pattern. All answered requests were resolved without clarification or semantic‑layer updates. For highly exploratory analytics—ad hoc queries that resist classification, novel join patterns or deeply nested aggregations—an unconstrained system will achieve higher raw coverage. AEGIS is not designed for that use case. It is designed for the majority of institutional reporting needs that fit within a bounded, maintainable semantic model, delivering guaranteed safety and semantic fidelity within that scope. Organizations choosing AEGIS trade unrestricted expressiveness for auditability, safety and reuse. Whether that trade is appropriate depends on the reporting environment. In institutional contexts where data privacy, regulatory compliance and reporting consistency matter, the tradeoff favors rigid system guarantees. In open‑ended data‑science exploration it does not.

\subsection{Widget Persistence as a Design Goal}
My evaluation confirms that widget reuse is a high‑value behaviour that does not emerge spontaneously in conventional reporting tools. In my baseline simulations equivalent reports were recreated multiple times because the system lacked any notion of persistence or similarity. AEGIS’s similarity‑based widget retrieval automatically surfaces existing widgets when available, demonstrating that persistence and reuse should be primary design goals in institutional NL reporting systems.

\subsection{Visualization Selection Quality}
In informal demonstrations of AEGIS I found that users rarely disagreed with the system’s default chart selection. Only a handful of instances prompted a different chart preference, typically for small‑cardinality categorical comparisons where a pie chart was preferred. These observations suggest that the visualization policy covers most user preferences and that the override‑with‑explanation behaviour (Section 4.9) is perceived as informative rather than paternalistic.

\section{Limitations and Future Work\label{sec:limitations}}
\paragraph{Benchmark Selection.} Standard NL2SQL datasets (e.g., Spider) test schema-linking by generating complex structural queries across arbitrary databases. AEGIS solves a different problem: reliable enterprise reporting over a fixed schema. The custom 100-query benchmark is necessary because standard benchmarks do not evaluate adversarial safety (SQL injection attempts) or strict adherence to pre-defined business logic (e.g., internal KPI formulas).

\paragraph{Architectural Overhead.} Adding a semantic mapper and safe query compiler introduces execution steps absent in direct LLM-to-SQL systems. However, this overhead is mathematically negligible. The compiler module (AST generation, validation, permission rewriting) executes in <10 milliseconds, representing less than 1\% of the total request latency when compared to typical LLM API inference times (1-3 seconds).

\paragraph{Semantic Layer Scalability.} The prototype injects the entire semantic vocabulary into the system prompt. While seemingly limited, modern 128k context windows can hold approximately 2,500 distinct metric and dimension definitions. Since most enterprise deployments expose fewer than 500 core reporting concepts, context limits are not a practical constraint. Furthermore, using compact data structures and short reference aliases heavily compresses the required tokens. Future work for massive-scale deployments (tens of thousands of concepts) could incorporate Retrieval-Augmented Generation (RAG) to dynamically inject only relevant semantic subsets.

\paragraph{Database Agnosticism.} The compiler implementation currently generates MySQL syntax. Because the architecture decouples intent extraction from syntax generation, supporting PostgreSQL or SQL Server requires only extending the compiler module; the LLM prompts and semantic mapping logic remain completely unchanged.

\paragraph{Storage Persistence.} The prototype uses JSON flat files to store widgets, prioritizing open-source portability and research reproducibility without requiring complex database installations. The widget registry interface is designed to seamlessly swap to relational databases (e.g., PostgreSQL) for commercial deployments.

\paragraph{Semantic layer maintenance.} The semantic layer requires ongoing maintenance as reporting needs evolve. While the one‑time construction effort for the e-commerce domain was approximately 40 person‑hours, ongoing maintenance adds overhead that small organizations may find burdensome. Future work will explore semi‑automated semantic layer extension from historical query logs and administrator feedback.

\paragraph{Template coverage.} The 2.6\% of requests outside the template library are not edge cases of adversarial prompting; they are legitimate analytical needs (cohort comparisons, multi‑level groupings) that require richer patterns. Extending the template library is straightforward but increases the validation and maintenance surface. I plan to add 4–6 additional patterns in a follow‑up study.

\paragraph{Parameter coercion edge case.} The residual 0\% unsafe query rate arises from implicit type coercion when a bound parameter value matches a row predicate unexpectedly. This can be addressed by adding domain validation of parameter values against their semantic‑layer datatype definition before binding.

\paragraph{Generalization beyond the e-commerce domain.} AEGIS is evaluated on an e-commerce retail domain. The architecture is domain‑agnostic, but each deployment requires a new semantic layer. Future work will evaluate AEGIS in financial reporting and healthcare administrative domains.

\paragraph{LLM cost.} Each request requires one LLM API call (two if a repair is needed). At current API pricing, this adds marginal latency (~800 ms) and cost per request. For high‑volume reporting environments, semantic caching of intent objects for identical requests would reduce both.

\paragraph{Multi‑turn conversation.} AEGIS currently treats each request independently. Users who issue follow‑up queries (``now filter that to just Electronics category'') must re‑specify the full request. Contextual carryover across turns is planned as the next major feature.

\section{Conclusion}
AEGIS presents a governed architecture for converting natural‑language reporting requests into safe, reusable dashboard widgets over relational databases. The core contribution is architectural: by confining the LLM to structured intent extraction and entrusting all execution to deterministic, bounded components, AEGIS eliminates the structural source of unsafe SQL generation while preserving the natural‑language usability that makes LLM‑assisted interfaces valuable. A formative study of real institutional reporting requests grounds the design, showing that eleven analytics primitives cover 98.2\% of realistic needs and that 61\% of requests are recurrences—findings that directly motivate the template library and widget reuse model. My quantitative evaluation on a 100‑query benchmark shows that AEGIS reduces unsafe SQL generation from 100\% (baseline) to 0\%, achieves 100\% intent accuracy and 100\% semantic correctness, and maintains execution validity at 100\%. 

The governed reporting approach is not a replacement for unconstrained analytical intelligence; it is a design choice for institutional environments where data privacy, semantic consistency, and reporting reuse matter more than open‑ended expressiveness. I argue that this environment describes the majority of real‑world business intelligence deployments, and that AEGIS demonstrates a practical path to making LLM‑assisted analytics trustworthy in production.

\section*{Declarations}
\begin{description}
\item[Funding:] No funding was received for this study.
\item[Conflict of Interest:] The author declares no conflict of interest.
\item[Data Availability:] The benchmark dataset, semantic layer configuration files, and evaluation scripts will be released publicly upon paper acceptance.
\item[Code Availability:] The AEGIS prototype implementation will be released as open‑source software upon paper acceptance.
\end{description}

\begin{thebibliography}{99}
\bibitem{affolter2019survey} Affolter, K., Stockinger, K., \& Bernstein, A. (2019). A comparative survey of recent natural language interfaces for databases. \emph{The VLDB Journal}, \emph{28}, 793--819. \url{https://doi.org/10.1007/s00778-019-00567-8}
\bibitem{deng2023dashbot} Deng, D., Wu, A., Qu, H., \& Wu, Y. (2023). DashBot: Insight-driven dashboard generation based on deep reinforcement learning. \emph{IEEE Transactions on Visualization and Computer Graphics}, \emph{29}(1), 690--700. \url{https://doi.org/10.1109/TVCG.2022.3209468}
\bibitem{gao2015datatone} Gao, T., Dontcheva, M., Adar, E., Liu, Z., \& Karahalios, K.~G. (2015). DataTone: Managing ambiguity in natural language interfaces for data visualization. \emph{Proceedings of the 28th Annual ACM Symposium on User Interface Software \& Technology (UIST)}, 489--500. \url{https://doi.org/10.1145/2807442.2807478}
\bibitem{lehmann2022building} Lehmann, C., Kehlbeck, R., Fekete, J.-D., \& Deussen, O. (2022). Building natural language interfaces for databases in practice. \emph{Proceedings of the 34th International Conference on Scientific and Statistical Database Management (SSDBM)}, Article~20. \url{https://doi.org/10.1145/3538712.3538744}
\bibitem{li2014nalir} Li, F., \& Jagadish, H.~V. (2014). Constructing an interactive natural language interface for relational databases. \emph{Proceedings of the VLDB Endowment}, \emph{8}(1), 73--84. \url{https://doi.org/10.14778/2735461.2735468}
\bibitem{li2023bigbench} Li, J., Hui, B., Qu, G., Yang, J., Li, B., Li, B., Wang, B., Qin, B., Geng, R., Huo, N., Zhou, X., Ma, C., Li, G., Chang, K.~C.-C., Qin, F., Cheng, R., \& Li, Y. (2023). Can large language models serve as a database interface? A big bench for large-scale database grounded text-to-SQLs. \emph{Advances in Neural Information Processing Systems (NeurIPS)}, \emph{36}.
\bibitem{liu2026systematic} Liu, M., Yang, H., Zhang, H., Zhang, Y., Wang, Y., \& Chen, Y. (2026). A systematic review of natural language interfaces for databases. \emph{Frontiers of Computer Science}, \emph{20}, 2011623. \url{https://doi.org/10.1007/s11704-025-50592-w}
\bibitem{luo2021nl2vis} Luo, Y., Tang, N., Li, G., Tang, J., Chai, C., \& Qin, X. (2021). Synthesizing natural language to visualization (NL2VIS) benchmarks from NL2SQL benchmarks. \emph{Proceedings of the 2021 International Conference on Management of Data (SIGMOD)}, 1235--1247. \url{https://doi.org/10.1145/3448016.3457259}
\bibitem{narechania2021nl4dv} Narechania, A., Srinivasan, A., \& Stasko, J. (2021). nl4dv: A toolkit for generating analytic specifications for data visualization from natural language queries. \emph{IEEE Transactions on Visualization and Computer Graphics}, \emph{27}(2), 369--379. \url{https://doi.org/10.1109/TVCG.2020.3030378}
\bibitem{openai2024structured} OpenAI. (2024). \emph{Introducing structured outputs in the API}. \url{https://openai.com/index/introducing-structured-outputs-in-the-api/}
\bibitem{postgresql} PostgreSQL Global Development Group. (2026). \emph{PostgreSQL documentation: CREATE POLICY}. \url{https://www.postgresql.org/docs/current/sql-createpolicy.html}
\bibitem{scholak2021picard} Scholak, T., Schucher, N., \& Bahdanau, D. (2021). PICARD: Parsing incrementally for constrained auto-regressive decoding from language models. \emph{Proceedings of the 2021 Conference on Empirical Methods in Natural Language Processing (EMNLP)}, 9895--9901. \url{https://doi.org/10.18653/v1/2021.emnlp-main.779}
\bibitem{setlur2016eviza} Setlur, V., Battersby, S.~E., Tory, M., Gossweiler, R., \& Chang, A.~X. (2016). Eviza: A natural language interface for visual analysis. \emph{Proceedings of the 29th Annual ACM Symposium on User Interface Software \& Technology (UIST)}, 365--377. \url{https://doi.org/10.1145/2984511.2984588}
\bibitem{shalaan2025gsql} Shalaan, H.~S., Soliman, T.~H.~A., \& AbdelAziz, A.~M. (2025). G-SQL: A schema-aware and rule-guided approach for robust natural language to SQL translation. \emph{IEEE Access}, \emph{13}, 158520--158534. \url{https://doi.org/10.1109/ACCESS.2025.3607879}
\bibitem{shailesh2025convbi} Shailesh, G.~N., Pavithran, M., Venkat, R.~H.~A., \& Kaliappan, P. (2025). Conversational BI: Natural language interface to business dashboards. \emph{International Journal of Engineering Research \& Technology}, \emph{14}(12). \url{https://doi.org/10.17577/IJERTV14IS120229}
\bibitem{su2026trisql} Su, X., Gu, Y., Wang, P., Gu, W., Qi, L., \& He, J. (2026). A robust natural language text-to-SQL generation framework with dynamic strategies based on large language models. \emph{Scientific Reports}, \emph{16}, Article~7892. \url{https://doi.org/10.1038/s41598-026-39128-9}
\bibitem{wang2020ratsql} Wang, B., Shin, R., Liu, X., Polozov, O., \& Richardson, M. (2020). RAT-SQL: Relation-aware schema encoding and linking for text-to-SQL parsers. \emph{Proceedings of the 58th Annual Meeting of the Association for Computational Linguistics (ACL)}, 7567--7578. \url{https://doi.org/10.18653/v1/2020.acl-main.677}
\bibitem{yu2018spider} Yu, T., Zhang, R., Er, H.~Y., Li, S., Xue, E., Pang, B., \ldots\ Radev, D. (2018). Spider: A large-scale human-labeled dataset for complex and cross-domain semantic parsing and text-to-SQL task. \emph{Proceedings of the 2018 Conference on Empirical Methods in Natural Language Processing (EMNLP)}, 3911--3921. \url{https://doi.org/10.18653/v1/D18-1425}
\bibitem{yu2019sparc} Yu, T., Zhang, R., Yasunaga, M., Tan, Y.~C., Lin, X.~V., Li, S., \ldots\ Xiong, C. (2019a). SParC: Cross-domain semantic parsing in context. \emph{Proceedings of the 57th Annual Meeting of the Association for Computational Linguistics (ACL)}, 4511--4523. \url{https://doi.org/10.18653/v1/P19-1443}
\bibitem{yu2019cosql} Yu, T., Zhang, R., Er, H., Li, S., Xue, E., Pang, B., \ldots\ Radev, D. (2019b). CoSQL: A conversational text-to-SQL challenge towards cross-domain natural language interfaces to databases. \emph{Proceedings of the 2019 Conference on Empirical Methods in Natural Language Processing (EMNLP)}, 1962--1979. \url{https://doi.org/10.18653/v1/D19-1204}
\bibitem{zhong2018seq2sql} Zhong, V., Xiong, C., \& Socher, R. (2018). Seq2SQL: Generating structured queries from natural language using reinforcement learning. \emph{Proceedings of the International Conference on Learning Representations (ICLR)}.
\end{thebibliography}

\end{document}
